\chapter{Commodity Options and Structured Hedges}\label{ch:m3:commodity-options-and-structured-hedges}

Every year Mexico's finance ministry buys put options on the oil it expects to export. In
December 2009, when prices had collapsed, its options paid almost 5.1 billion dollars. In 2015 they paid 6.4 billion, and 2.7 billion
the year after. It is one of the largest regular hedges in any market, and it is built from
the instruments of this chapter: options on the average of a price over a period, bought from
banks that hedge them in the futures market. Commodity producers, consumers and governments
hedge not a price on a day but a revenue or a cost over months, and their instruments are
shaped accordingly. This chapter covers options on commodity futures, average-price and
\omterm{def:m3:commodity-options-and-structured-hedges:cso}{calendar-spread options}, the \omterm{def:m3:commodity-options-and-structured-hedges:swing}{swing contracts} of gas and power, the \omterm{def:m3:commodity-options-and-structured-hedges:collar}{collars} that producers sell
to finance their protection, and the \omterm{def:m3:commodity-options-and-structured-hedges:programme}{hedging programmes} that combine them.

\section{Options on commodity futures}

\begin{definition}[Futures option]\label{def:m3:commodity-options-and-structured-hedges:fo}
A \emph{futures option}\index{futures option} is an option whose exercise delivers a position in a
futures contract (long for a call, short for a put) at the strike, the difference being settled
through the futures' variation margin; it usually expires shortly before the future it is written on.
\end{definition}

Because the underlying is a future, which costs nothing to hold, a \omterm{def:m3:commodity-options-and-structured-hedges:fo}{futures option} is priced on the
forward, not on a spot plus carry: Black's formula for options on futures, which One Quant Book 5
derives, uses the futures price and discounts the payoff to the option's expiry. For a future at \$60,
a one-year put struck at \$55 with a volatility of 35\% and a rate of 4\% is worth \$5.51 a barrel. NYMEX's
listed WTI option is American and expires at the close three business days before its future; its WTI
\omterm{def:m3:commodity-options-and-structured-hedges:apo}{average-price option} is European, cash-settled on the month's average settlement of the nearby future times
1{,}000 barrels, and expires on the month's last business day.

Commodity options differ from equity options in the shape of their volatility smile (One Quant Book 1,
chapter 25). Where supply shocks push prices up (a refinery outage, a frost, a war in a producing
region), out-of-the-money calls can be as dear as puts, or dearer; where a glut is the fear, puts are.
And near-dated options carry the higher volatility of near-dated futures, the \omterm{def:m3:forward-curves-storage-and-convenience-yield:samuelson}{Samuelson effect} of
\cref{ch:m3:forward-curves-storage-and-convenience-yield}.

\section{Average-price and calendar-spread options}

A producer that sells oil every day of a month is exposed to the month's average price, not to its last
day's. Its option is on that average.

\begin{definition}[Average-price option]\label{def:m3:commodity-options-and-structured-hedges:apo}
An \emph{average-price option}\index{average-price option} (APO) pays at expiry on the arithmetic
average of a reference price over a period (a month, a quarter, a year) against its strike: a put
pays $\max(K - \bar P, 0)$ per unit. It is the product commodity markets trade under that name; One
Quant Book 5, chapter 16, prices the general Asian option.
\end{definition}

An average moves less than its last observation: with daily fixings, a lognormal price's average over
the option's life has about $1/\sqrt 3$ of the price's volatility. The average-price put of the example
above, on twelve monthly fixings, is worth \$2.66 a barrel by simulation, less than half the vanilla put's
\$5.51 (\cref{fig:m3:commodity-options-and-structured-hedges:dist}).

\begin{omfigure}
\begin{tikzpicture}
\begin{axis}[width=11cm,height=4.8cm,ybar,bar width=4pt,
  xlabel={\$ per barrel},ylabel={\% of paths},
  tick label style={font=\small},label style={font=\small},
  xmin=10,xmax=150,ymin=0,ymax=22,grid=major,grid style={gray!25},
  legend columns=2,legend style={font=\footnotesize,at={(0.5,-0.33)},anchor=north,draw=none,fill=none,/tikz/every even column/.append style={column sep=8pt}},
  table/col sep=comma]
\addplot[fill=omDef!55,draw=omDef] table[x=mid,y=terminal]{figdata/markets-3/12-commodity-options-and-structured-hedges/dist.csv};
\addplot[fill=omThm!55,draw=omThm] table[x=mid,y=average]{figdata/markets-3/12-commodity-options-and-structured-hedges/dist.csv};
\legend{futures price in one year,average of twelve monthly fixings}
\end{axis}
\end{tikzpicture}

\omcaption{The futures price at the end of a year and the average of its twelve monthly fixings, both
starting at \$60 with 35\% volatility, 200{,}000 simulated paths in bins of \$5: the average is far less
dispersed, so options on it are cheaper. Illustrative; data: the chapter's tutorial.}\label{fig:m3:commodity-options-and-structured-hedges:dist}
\end{omfigure}

\begin{definition}[Calendar-spread option]\label{def:m3:commodity-options-and-structured-hedges:cso}
A \emph{calendar-spread option}\index{calendar-spread option} (CSO) is an option on the difference between
the prices of two delivery months of the same commodity, for example a call paying $\max(F_{T_1} - F_{T_2}
- K, 0)$ on the first minus the second month.
\end{definition}

A storage owner is naturally long \omterm{def:m3:commodity-options-and-structured-hedges:cso}{calendar-spread options} (it profits when near prices rise above far
ones by more than the cost of carry), and a refiner or a shipper is long or short them according to its
timing. Their prices depend on the correlation between the two months, which is high but falls when the
curve's shape changes; One Quant Book 6, chapter 16, models spread options.

\section{Swing contracts}

Gas and power are bought under contracts that give the buyer flexibility in how much it takes, because
its own demand varies with the weather.

\begin{definition}[Swing contract, take-or-pay clause]\label{def:m3:commodity-options-and-structured-hedges:swing}
A \emph{swing contract}\index{swing contract} lets its holder choose, each day, a quantity between a
minimum and a maximum around a daily contract quantity, at a fixed or indexed price, subject to limits on
the total taken over the year. A \emph{take-or-pay clause}\index{take-or-pay clause} obliges the buyer to
pay for a minimum annual quantity whether or not it takes it.
\end{definition}

\begin{omfigure}
\centering
\begin{tikzpicture}[font=\small]
\begin{axis}[width=11cm,height=4.4cm,
  xlabel={day of the contract year (stylised)},ylabel={daily quantity},
  tick label style={font=\small},label style={font=\small},
  xmin=0,xmax=365,ymin=0,ymax=150,xtick=\empty,ytick={80,100,120},yticklabels={min,DCQ,max}]
\addplot[omDef,thick,dashed,domain=0:365] {120};
\addplot[omDef,thick,dashed,domain=0:365] {80};
\addplot[gray,thin,domain=0:365] {100};
\addplot[omThm,very thick,const plot] coordinates {(0,118) (60,120) (100,95) (160,80) (240,82) (300,112) (365,112)};
\node[font=\footnotesize,anchor=south west] at (axis cs:5,122) {take more when the market price is high};
\node[font=\footnotesize,anchor=north west] at (axis cs:125,78) {take the minimum when it is low};
\end{axis}
\end{tikzpicture}

\omcaption{A \omterm{def:m3:commodity-options-and-structured-hedges:swing}{swing contract}, schematic: each day the holder nominates between a minimum and a maximum
around the daily contract quantity (DCQ); the year's total must stay within annual limits, and the
\omterm{def:m3:commodity-options-and-structured-hedges:swing}{take-or-pay clause} makes the lowest total a payment in any case. Stylised, no data.}\label{fig:m3:commodity-options-and-structured-hedges:swing}
\end{omfigure}

A \omterm{def:m3:commodity-options-and-structured-hedges:swing}{swing contract} is a strip of daily options to buy more or less at the contract price, linked by the
annual limits: taking the maximum today uses up flexibility that may be worth more later. Its value
therefore depends on the whole path of prices, and its optimal use is a dynamic programme like the
battery of \cref{ch:m3:power-markets-trading}; the model is One Quant Book 6's.

\begin{example}[Three days of swing]\label{ex:m3:commodity-options-and-structured-hedges:swing}
A buyer holds a three-day contract at \qty{40}{EUR/MWh} with a daily contract quantity of 100~MWh, a swing of
20~MWh either way, and a total that must equal 300~MWh. Market prices turn out to be 30, 50 and 40. Taking 80
on the first day (buying the missing 20 in the market at 30) and 120 on the second (selling the extra 20 at
50) is worth $20 \times (40 - 30) + 20 \times (50 - 40) = 400$ euros more than taking 100 every day. With
prices unknown in advance, the holder must decide each day, and the flexibility is worth less than this
perfect-foresight value.
\end{example}

\section{Collars and three-way collars}

Puts cost money. A producer that wants a floor but will not pay for it sells upside instead.

\begin{definition}[Collar, three-way collar]\label{def:m3:commodity-options-and-structured-hedges:collar}
A producer's \emph{collar}\index{collar} buys a put and sells a call, usually with strikes chosen so that
the premiums cancel (a zero-cost collar): its realised price is kept between the two strikes. A
\emph{three-way collar}\index{three-way collar} adds a sold put below the bought one: the producer is
protected only between the two put strikes, and the extra premium lets it raise the floor or the cap.
\end{definition}

\begin{omfigure}
\begin{tikzpicture}
\begin{axis}[width=11cm,height=5.2cm,
  xlabel={average price, \$/bbl},ylabel={realised price, \$/bbl},
  tick label style={font=\small},label style={font=\small},
  xmin=20,xmax=110,ymin=20,ymax=110,ytick={40,60,80,100},grid=major,grid style={gray!25},
  legend columns=2,legend style={font=\footnotesize,at={(0.5,-0.25)},anchor=north,draw=none,fill=none,/tikz/every even column/.append style={column sep=8pt}},
  table/col sep=comma]
\addplot[gray,thick] table[x=x,y=unhedged]{figdata/markets-3/12-commodity-options-and-structured-hedges/payoffs.csv};
\addplot[omDef,very thick] table[x=x,y=collar]{figdata/markets-3/12-commodity-options-and-structured-hedges/payoffs.csv};
\addplot[omThm,very thick,dashed] table[x=x,y=threeway]{figdata/markets-3/12-commodity-options-and-structured-hedges/payoffs.csv};
\addplot[omProp,thick,dotted] table[x=x,y=put]{figdata/markets-3/12-commodity-options-and-structured-hedges/payoffs.csv};
\legend{unhedged,collar 50/75,three-way 60/45/70.24,55 put net of premium}
\end{axis}
\end{tikzpicture}

\omcaption{A producer's realised price per barrel against the year's average price. The \omterm{def:m3:commodity-options-and-structured-hedges:collar}{collar} keeps it
between \$50 and \$75; the three-way protects from \$60 down to \$45 only, then follows the price down with
a \$15 cushion, and gives up everything above \$70.24; the bought put keeps all the upside at the cost of
its premium. Illustrative; data: the chapter's tutorial.}\label{fig:m3:commodity-options-and-structured-hedges:payoffs}
\end{omfigure}

\begin{example}[What each structure costs]\label{ex:m3:commodity-options-and-structured-hedges:costs}
With the future at \$60, 35\% volatility and a year of monthly averaging, the 50/75 \omterm{def:m3:commodity-options-and-structured-hedges:collar}{collar} costs
\$0.07 a barrel, close to zero; the 55 average-price put costs \$2.66. A three-way that buys the 60 put and
sells the 45 put costs the same \$2.66 if it also sells a call at \$70.24. A producer with a budget of \$50
a barrel would, unhedged, fall \$8.43 short of it in the worst 5\% of years; with either the \omterm{def:m3:commodity-options-and-structured-hedges:collar}{collar} or the
put its worst-5\% realised price stays above the budget.
\end{example}

The three-way was popular with oil producers because it looks cheap and protects against moderate falls.
Its weakness is the sold put: in a collapse the producer's protection stops at the lower strike just when
it is most needed.

\begin{omfigure}
\begin{tikzpicture}
\begin{axis}[width=11cm,height=4.8cm,
  xlabel={realised price over the year, \$/bbl},ylabel={cumulative probability},
  tick label style={font=\small},label style={font=\small},
  xmin=25,xmax=110,ymin=0,ymax=1,grid=major,grid style={gray!25},
  legend columns=3,legend style={font=\footnotesize,at={(0.5,-0.3)},anchor=north,draw=none,fill=none,/tikz/every even column/.append style={column sep=8pt}},
  table/col sep=comma]
\addplot[gray,thick] table[x=unhedged,y=p]{figdata/markets-3/12-commodity-options-and-structured-hedges/revcdf.csv};
\addplot[omDef,very thick] table[x=collar,y=p]{figdata/markets-3/12-commodity-options-and-structured-hedges/revcdf.csv};
\addplot[omProp,very thick,dashed] table[x=put,y=p]{figdata/markets-3/12-commodity-options-and-structured-hedges/revcdf.csv};
\legend{unhedged,collar 50/75,55 put net of premium}
\end{axis}
\end{tikzpicture}

\omcaption{Distribution of a producer's realised average price over a year, simulated from a \$60 future with
35\% volatility: unhedged, with the 50/75 \omterm{def:m3:commodity-options-and-structured-hedges:collar}{collar}, and with the 55 average-price put net of its \$2.66 premium.
The \omterm{def:m3:commodity-options-and-structured-hedges:collar}{collar} truncates both tails; the put truncates only the lower one, at a cost. Illustrative; data: the
chapter's tutorial.}\label{fig:m3:commodity-options-and-structured-hedges:revcdf}
\end{omfigure}

\section{Producer, consumer and sovereign hedging programmes}

\begin{definition}[Hedging programme]\label{def:m3:commodity-options-and-structured-hedges:programme}
A \emph{hedging programme}\index{hedging programme} is a standing policy, set by a board or a ministry,
fixing what share of future production or consumption is hedged, over what horizon, with which
instruments, and how the hedges are executed and reported.
\end{definition}

A producer hedges to protect its investment budget and its debt covenants; a consumer (an airline, a
utility, a food processor) to protect its margins; a government to protect its budget. They choose
differently: producers favour swaps and \omterm{def:m3:commodity-options-and-structured-hedges:collar}{collars} that cost nothing upfront, consumers calls or \omterm{def:m3:commodity-options-and-structured-hedges:collar}{collars}, and
governments puts, which cost a known premium and leave the upside to the treasury.

\begin{dated}{2026-09}{Mexico's oil hedge}\label{dat:m3:commodity-options-and-structured-hedges:mexico}
Mexico's government buys, each year, Asian (average-price) put options on its oil exports for the year
ahead, from banks. Reported payouts: almost USD~5.1 billion for 2009, USD~6.4 billion for 2015 and USD~2.7
billion for 2016. On the cost, in January 2020 the finance minister put it at about USD~1 billion a year and a
deputy minister at USD~1.2 billion on average over the programme's life, as reported by the press.
\end{dated}

\section{Tutorial: pricing an average-price put and a three-way}\label{tut:m3:commodity-options-and-structured-hedges:apo}

\begin{tutorial}
\textbf{Goal.} Price a vanilla put on a future and an average-price put, and solve the call strike of a
\omterm{def:m3:commodity-options-and-structured-hedges:collar}{three-way collar} that costs as much as the put. \textbf{End state:}
\cref{fig:m3:commodity-options-and-structured-hedges:dist,fig:m3:commodity-options-and-structured-hedges:payoffs}
and \cref{ex:m3:commodity-options-and-structured-hedges:costs}.
\begin{tutsteps}
\item \textbf{The legs.} A leg is a swap, put or call, bought or sold; structures are lists of legs.
\omcode{code/firm/hedgeprog/firm_hedgeprog.py}{28}{52}{Legs, collars and three-way collars.}\label{lst:m3:commodity-options-and-structured-hedges:legs}
\item \textbf{Simulate and value.} Monthly averages of a driftless lognormal future, then mean discounted
payoffs.
\omcode{code/firm/hedgeprog/firm_hedgeprog.py}{55}{68}{Averages by simulation and the value of a structure.}\label{lst:m3:commodity-options-and-structured-hedges:value}
\item \textbf{Run} \texttt{m3\_options.vanilla\_vs\_average()}, \texttt{sovereign()},
\texttt{producer\_car()} and \texttt{fig\_options.py}.
\end{tutsteps}
\textbf{What to change next.} Use daily instead of monthly fixings and compare with the $1/\sqrt3$ rule;
add a volatility skew (a higher volatility for low strikes) and see which structure becomes dearer.
\end{tutorial}

\section{Build: the hedging-programme evaluator}\label{bld:m3:commodity-options-and-structured-hedges:hedgeprog}

\begin{build}
\bfield{Purpose} The miniature firm structures hedges for producers and consumers: it must value their
structures, show their payoffs, and measure what each does to the client's worst years.
\bfield{Interface} \texttt{black76(f, k, t, sigma, r, call)}; \texttt{Leg(kind, strike, side).payoff};
\texttt{collar}, \texttt{three\_way}; \texttt{simulate\_averages(f0, sigma, t, steps, n, seed)};
\texttt{value(legs, settlements, r, t)}; \texttt{hedged\_revenue(price, legs, premium)};
\texttt{cash\_flow\_at\_risk(revenue, budget, q)}.
\bfield{Rules} Payoffs per barrel on the settlement price (the average for average-price legs); simulation
seeds fixed; the future is a martingale under the pricing measure; cash flow at risk is the shortfall of
the lower quantile below the budget.
\bfield{Acceptance tests} \texttt{code/firm/hedgeprog/tests/}: put--call parity for Black's formula; the
average-price put cheaper than the vanilla and the simulated average centred on the forward; a \omterm{def:m3:commodity-options-and-structured-hedges:collar}{collar}
bounding revenue; a three-way's capped protection.
\bfield{Stretch} Volatility smiles by strike; quarterly and annual averaging periods; the producer's
volumes as uncertain as its prices.
\end{build}

\begin{omsources}
\item MoneyWeek, ``Mexico's lucrative Hacienda hedge'', 21 April 2017.
\item Jain Family Institute, \emph{Mexico's Petroleum Hedging Program}, 2023.
\item Alto Nivel (with Reuters), ``¿Qué son las coberturas petroleras y cuánto le cuestan a México?'', 10 January 2020.
\item F.~Black, ``The pricing of commodity contracts'', \emph{Journal of Financial Economics}, 1976.
\item CME Group, WTI average price options; NYMEX Rulebook, chapters 310 (Light Sweet Crude Oil Option) and 341 (WTI Average
Price Option) (Internet Archive copies).
\end{omsources}

\section{Exercises}

\begin{exercise}[$\star$]\label{exo:m3:commodity-options-and-structured-hedges:1}
A producer sells a swap at \$62 for a year of monthly averages. The year averages \$55. What does the swap
pay, and what is the producer's realised price?
\end{exercise}

\begin{exercise}[$\star$]\label{exo:m3:commodity-options-and-structured-hedges:2}
Why is an \omterm{def:m3:commodity-options-and-structured-hedges:apo}{average-price option} cheaper than the vanilla option with the same strike and expiry?
\end{exercise}

\begin{exercise}[$\star$]\label{exo:m3:commodity-options-and-structured-hedges:3}
Give the realised price per barrel of the 50/75 \omterm{def:m3:commodity-options-and-structured-hedges:collar}{collar} if the year averages \$40, \$60 and \$90.
\end{exercise}

\begin{exercise}[$\star\star$]\label{exo:m3:commodity-options-and-structured-hedges:4}
For the three-way 60/45/70.24, give the realised price if the year averages \$30, \$50, \$65 and \$90.
\end{exercise}

\begin{exercise}[$\star\star$]\label{exo:m3:commodity-options-and-structured-hedges:5}
With twelve monthly fixings, the volatility of the average's logarithm in the simulation is 21.3\%. Compare
it with $35\%/\sqrt3$ and explain the difference.
\end{exercise}

\begin{exercise}[$\star\star$]\label{exo:m3:commodity-options-and-structured-hedges:6}
Why does a storage owner hold, in effect, \omterm{def:m3:commodity-options-and-structured-hedges:cso}{calendar-spread options}?
\end{exercise}

\begin{exercise}[$\star\star\star$]\label{exo:m3:commodity-options-and-structured-hedges:7}
\emph{Coding.} With \texttt{black76}, verify put--call parity for the \omterm{def:m3:commodity-options-and-structured-hedges:fo}{futures option} of the text, and give
the call's value.
\end{exercise}

\begin{exercise}[$\star\star\star$]\label{exo:m3:commodity-options-and-structured-hedges:8}
\emph{Find the flaw.} ``Our \omterm{def:m3:commodity-options-and-structured-hedges:collar}{three-way collar} costs nothing and protects us from \$60 down.''
\end{exercise}

\section{Problem: The Finance Ministry's Put}

\begin{problem}[{Weekend problem --- a sovereign's annual oil put}]\label{pb:m3:commodity-options-and-structured-hedges:1}
A finance ministry expects to export 250 million barrels next year and budgets them at the forward, \$60. It
can buy an average-price put at \$55, or spend the same premium on a three-way: buy the 60 put, sell the 45
put, and sell a call. Volatility is 35\%, the rate 4\%, and averaging monthly (all illustrative).

\medskip\noindent\textbf{Part I --- The put.}
\begin{enumerate}
\item What is the put's premium per barrel?
\item What does the programme cost, in dollars?
\item What share of the budgeted export revenue is that?
\item What does the put pay if the year averages \$45?
\item And if it averages \$30?
\end{enumerate}

\medskip\noindent\textbf{Part II --- The three-way.}
\begin{enumerate}[resume]
\item At what strike must the call be sold for the same premium?
\item What does the three-way pay at an average of \$45?
\item And at \$30?
\item What does it cost the ministry if the year averages \$90?
\item Which structure does better in which years?
\end{enumerate}

\medskip\noindent\textbf{Part III --- Execution.}
\begin{enumerate}[resume]
\item Why buy \omterm{def:m3:commodity-options-and-structured-hedges:apo}{average-price options} rather than vanillas?
\item How do the selling banks hedge, and what does that do to the futures market?
\item Why might the ministry spread its purchases over weeks and keep them secret?
\item What does the ministry do if its exports fall short of 250 million barrels?
\item How should the premium be accounted for in the budget?
\end{enumerate}

\medskip\noindent\textbf{Part IV --- Judgement.}
\begin{enumerate}[resume]
\item Why might a government prefer puts to swaps?
\item What does the programme's history of payouts say about its value?
\item When is a three-way the wrong choice for a sovereign?
\item State the \emph{named result}: the cost of the put programme as a share of revenue, and its payouts in
a \$45 and a \$30 year against the three-way's.
\item In one sentence: what does a sovereign buy with a put?
\end{enumerate}
\end{problem}

\section{Interview questions}

\begin{interviewq}[$\star$ \iqroles{trader}]\label{iq:m3:commodity-options-and-structured-hedges:1}
What is an \omterm{def:m3:commodity-options-and-structured-hedges:apo}{average-price option}, and who needs it?
\end{interviewq}

\begin{interviewq}[$\star$ \iqroles{trader, bank}]\label{iq:m3:commodity-options-and-structured-hedges:2}
Explain a zero-cost \omterm{def:m3:commodity-options-and-structured-hedges:collar}{collar} to a producer's finance director.
\end{interviewq}

\begin{interviewq}[$\star\star$ \iqroles{researcher}]\label{iq:m3:commodity-options-and-structured-hedges:3}
Why is the effective volatility of an average about $1/\sqrt3$ of the price's?
\end{interviewq}

\begin{interviewq}[$\star\star$ \iqroles{bank, risk}]\label{iq:m3:commodity-options-and-structured-hedges:4}
You sold a large average-price put to a sovereign. How do you hedge it through the year?
\end{interviewq}

\begin{interviewq}[$\star\star$ \iqroles{trader}]\label{iq:m3:commodity-options-and-structured-hedges:5}
What is a \omterm{def:m3:commodity-options-and-structured-hedges:swing}{swing contract}, and why is it worth more than a fixed-volume contract?
\end{interviewq}

\begin{interviewq}[$\star\star\star$ \iqroles{developer, researcher}]\label{iq:m3:commodity-options-and-structured-hedges:6}
Design a tool that proposes and compares hedging structures for a producer, given its budget and covenants.
\end{interviewq}
